%% Jian Zhao — Master CV
%% AI Research Scientist | Applied Scientist | Machine Learning Engineer
%% Comprehensive reference — tailor role-specific CVs from this file.
%%
%% Compile from the repository root with:
%% lualatex --output-directory=cv-output cv-output/main_example-jian.tex
%% Positioning: Lead with Tsinghua M.S. + U.S. Ph.D. + Biomedical AI + Industry

\documentclass[11pt,a4paper,sans]{moderncv}
\moderncvstyle{banking}
\moderncvcolor{blue}

\renewcommand*{\firstnamestyle}[1]{{\fontsize{34}{36}\bfseries\upshape\color{color1}#1}}
\renewcommand*{\lastnamestyle}[1]{{\fontsize{34}{36}\bfseries\upshape\color{color1}#1}}
\renewcommand*{\sectionstyle}[1]{{\sectionfont\color{color1}#1}}

%%\usepackage[utf8]{inputenc}
\usepackage{microtype}

\usepackage{hyperref}
\hypersetup{
    colorlinks=true,
    linkcolor=blue,
    filecolor=magenta,
    urlcolor=blue,
    pdftitle={Jian Zhao - CV},
    pdfpagemode=FullScreen,
}
\usepackage[scale=0.80]{geometry}
\usepackage{needspace}
\usepackage{import}

\name{Jian}{Zhao}
\address{Philadelphia, PA}{}{}
\phone[mobile]{302-359-7657}
\email{zhaojianphd@gmail.com}
\extrainfo{\href{https://linkedin.com/in/jian-zhao-12800616a}{LinkedIn}}

\begin{document}

\makecvtitle

% ============================================================
%     PROFILE STATEMENT
% ============================================================

\vspace{6pt}
\small{AI researcher and software engineer with Tsinghua University graduate training, U.S. master's and Ph.D. research, and 10+ years across biomedical AI, graph learning, computer vision, LLM applications, HPC, and enterprise systems. Combines rigorous interdisciplinary research in healthcare and life sciences with production software engineering experience, multi-institution collaboration (Princeton, UCLA, Lawrence Livermore), and hands-on implementation in Python and PyTorch. H-1B transfer available.}

% ============================================================
%     CORE COMPETENCIES
% ============================================================

\section{Core Competencies}
\vspace{1pt}
\begin{itemize}

\item \textbf{Generative AI \& LLMs}: LLaMA, ChatGPT API, Hugging Face Transformers, PEFT (LoRA/QLoRA), Retrieval-Augmented Generation (RAG), prompt engineering, healthcare conversational AI; Milvus, Pinecone, Neo4j for vector and graph retrieval

\item \textbf{Deep Learning \& Graph ML}: PyTorch, TensorFlow, Scikit-learn, PyTorch Geometric; CNNs, Graph Neural Networks (GNN), temporal GNN, representation learning; biomedical cell segmentation and tracking, computer vision, object detection

\item \textbf{Healthcare AI \& Biomedical Research}: Clinical NLP, EHR data pipelines, public health analytics; biomedical image analysis (microscopy, cell tracking, segmentation); XGBoost, Random Forest, SVM for clinical ML; EMMA healthcare assistant (LLM/RAG, postpartum-depression support)

\item \textbf{HPC \& Scientific Computing}: MPI, MPI I/O, OpenMP, GPU computing; OSIRIS and GROMACS for plasma and molecular simulation; HPC cluster administration (320 CPU / 512 GB); LUSTRE, Hadoop distributed storage

\item \textbf{Engineering \& Systems}: Python (expert), Java (advanced), C++, C\#, MATLAB, R, SQL; PostgreSQL, Oracle, MySQL; Docker, Git, Linux; Apache Tomcat, PHP, WordPress; enterprise software architecture and delivery

\end{itemize}

% ============================================================
%     PROFESSIONAL AND RESEARCH EXPERIENCE
% ============================================================

\section{Professional and Research Experience}
\vspace{3pt}
\begin{itemize}

\needspace{5\baselineskip}
\item{\cventry{Oct 2023--Apr 2025}{Assistant Research Scientist}{Temple University, College of Public Health}{Philadelphia, PA}{}{\vspace{1pt}
\begin{itemize}
    \item Developed EMMA, an LLM-based conversational assistant for postpartum-depression support using LLaMA/ChatGPT, RAG, PEFT, Milvus, Neo4j, Flowise, Gradio, and Streamlit to match users with locally relevant mental-health resources
    \item Applied Random Forest, XGBoost, and SVM to ambulatory-visit and public-health datasets to identify monitoring gaps in blood-pressure and A1C tracking and support hypertension-prediction research
    \item Designed data and AI workflows spanning electronic health records, environmental sensors, and public-health repositories; contributed research results, manuscripts, grant proposals, and technical documentation
    \item Collaborated with faculty, clinicians, research staff, and community partners; co-supervised undergraduate and graduate researchers alongside the principal investigator
\end{itemize}}}

\vspace{3pt}

\needspace{5\baselineskip}
\item{\cventry{Aug 2019--Present}{PhD Research Assistant}{Delaware State University}{Dover, DE}{}{\vspace{1pt}
\begin{itemize}
    \item Researching biomedical cell segmentation, tracking, and quantification using mathematical algorithms, image processing, CNNs, GNNs, and LLM-assisted parameter tuning for microscopy applications
    \item Developed a rapid-tuning approach for CSTQ parameters to support reproducible biomedical image-analysis workflows and disease-research applications
    \item Conducting large-scale nano-bio spectroscopy and plasma simulations using MPI, MPI I/O, OpenMP, GPU computing, OSIRIS, and GROMACS in multi-institution collaborations with Princeton University, UCLA, and Lawrence Livermore National Laboratory
\end{itemize}}}

\vspace{3pt}

\needspace{5\baselineskip}
\item{\cventry{Aug 2016--Jul 2019}{Graduate Research Assistant}{Delaware State University}{Dover, DE}{}{\vspace{1pt}
\begin{itemize}
    \item Built and administered an HPC research cluster (1 head node, 8 compute nodes, 320 CPUs, 512 GB RAM) for nano-bio spectroscopy simulations and data-intensive experiments
    \item Completed machine-learning projects in pattern recognition and dimensionality reduction (PCA, LDA, LLE) using TensorFlow and Scikit-learn
    \item Investigated wireless insulin-pump security vulnerabilities and proposed a visible-light-communication access-control scheme using IoT prototypes
\end{itemize}}}

\vspace{3pt}

\needspace{5\baselineskip}
\item{\cventry{Aug 2007--Jul 2016}{Senior Software Engineer}{PetroChina Planning \& Engineering Institute}{Beijing, China}{}{\vspace{1pt}
\begin{itemize}
    \item Developed BP neural-network and decision-tree models for natural-gas pipeline energy-consumption forecasting, achieving approximately 90\% accuracy for strategic planning analyses
    \item Built and validated an M/M/C queuing model to optimize compressed-natural-gas vehicle allocation across varying fleet sizes
    \item Managed project delivery and resolved technical issues across analytics and enterprise software initiatives using MATLAB, Java, SQL, Weka, and statistical modeling
\end{itemize}}}

\vspace{3pt}

\needspace{5\baselineskip}
\item{\cventry{Aug 2005--Jul 2007}{Project Manager}{Chinese Academy of Sciences}{Beijing, China}{}{\vspace{1pt}
\begin{itemize}
    \item Led a five-person team on a Sino-German national forest-pest-control information-system project, delivering a scalable B/S application for the National Forestry program
    \item Modernized the Forest News application from ASP to JSP and migrated the database from Access to Oracle 9i using Java, Spring, Hibernate, Struts, JUnit, and GIS technologies
\end{itemize}}}

\end{itemize}

% ============================================================
%     SELECTED RESEARCH AND TECHNICAL PROJECTS
% ============================================================

\section{Selected Research and Technical Projects}
\vspace{1pt}
\begin{itemize}
\item \textbf{EMMA Healthcare Assistant}: LLM and RAG application for postpartum-depression support using LLaMA, ChatGPT, PEFT, Milvus, Neo4j, Hugging Face, Flowise, Gradio, and Streamlit.
\item \textbf{Biomedical Cell Tracking}: CNN/GNN and image-processing pipeline for cell segmentation, tracking, and quantification including watershed methods and diffusion-tensor imaging concepts.
\item \textbf{Ambulatory Health Analytics}: ML modeling of patient visits to identify blood-pressure and A1C monitoring gaps; hypertension-prediction research using ensemble methods.
\item \textbf{HPC Nano-Bio Simulation}: Parallel spectroscopy and plasma simulation using MPI, OpenMP, GPU computing, OSIRIS, GROMACS, and distributed storage (LUSTRE, Hadoop).
\end{itemize}

% ============================================================
%     EDUCATION
% ============================================================

\section{Education}
\vspace{1pt}
\begin{itemize}

\needspace{5\baselineskip}
\item{\cventry{Aug 2019--Expected 2026}{Ph.D. Candidate, Applied Mathematics and Mathematical Physics}{Delaware State University}{Dover, DE}{}{\vspace{1pt}
In progress; expected completion 2026. GPA: 3.89. Research: biomedical AI, graph neural networks, HPC simulation.
}}

\vspace{3pt}

\item{\cventry{Aug 2016--Jul 2019}{M.S., Computer and Information Sciences}{Delaware State University}{Dover, DE}{}{\vspace{1pt}
GPA: 3.9. Machine learning, pattern recognition, HPC cluster administration.
}}

\vspace{3pt}

\needspace{5\baselineskip}
\item{\cventry{2003--2005}{M.E., Software Engineering}{Tsinghua University}{Beijing, China}{}{\vspace{1pt}
GPA: 3.3. Software architecture, enterprise systems, project management.
}}

\vspace{3pt}

\item{\cventry{1998--2002}{B.E., Ocean Engineering; Minor in Computer Science}{Dalian University of Technology}{Dalian, China}{}{\vspace{1pt}
Engineering fundamentals, applied mathematics, computer science.
}}

\end{itemize}

% ============================================================
%     LANGUAGES
% ============================================================

\section{Languages}
\vspace{1pt}
\begin{itemize}
\item English (professional working proficiency), Chinese/Mandarin (native).
\end{itemize}

% ============================================================
%     PUBLICATIONS
% ============================================================

\section{Publications}
\vspace{1pt}
\begin{itemize}
\item Manuscripts contributed during Temple University research (2023--2025) — titles and venues to be confirmed.
\item Ph.D. research outputs in biomedical cell tracking and image analysis — titles and venues to be confirmed.
\end{itemize}

% ============================================================
%     ADDITIONAL
% ============================================================

\section{Additional}
\vspace{1pt}
\begin{itemize}
\item{\textbf{Work Authorization:} H-1B transfer available.}
\item{\textbf{Research Positioning:} Biomedical AI, graph machine learning, computer vision, LLM applications, and high-performance computing.}
\end{itemize}

% ============================================================
%     REFERENCES
% ============================================================

\section{References}
\vspace{1pt}
\begin{itemize}
\item \textbf{Dr. Sokratis Makrogiannis} --- Associate Professor in Mathematics, Delaware State University\\
\href{mailto:smakrogiannis@desu.edu}{smakrogiannis@desu.edu}

\item \textbf{Dr. Xiali Hei} --- Alfred and Helen M. Lamson Endowed Professor in Computing \& Informatics, University of Louisiana at Lafayette\\
\href{mailto:c00404592@louisiana.edu}{c00404592@louisiana.edu}

\item \textbf{Dr. Kam Kong} --- Associate Professor in Computer and Information Sciences, Delaware State University\\
\href{mailto:kkong@desu.edu}{kkong@desu.edu}
\end{itemize}

\end{document}
